\chapter{Implementation Status}

This appendix is a snapshot of built and verified behavior. Source, tests, and
repository history remain authoritative.

\section{Implementation phases}

\begin{longtable}{|p{1.1cm}|p{7.4cm}|p{3cm}|}
\hline
\textbf{Phase} & \textbf{Description} & \textbf{Status} \\
\hline
1 & Capability table, rights masks, object teardown & Done \\
2 & Endpoint IPC v1: scalar send/call/receive/reply, connection delegation,
    reply tokens & Done \\
3 & Userspace name/service manager, bootstrap delivery, ELF loader,
    supervisor, generation tracking & Done \\
4 & First-class memory objects: allocate, map, unmap, close, ownership
    accounting & Done \\
5 & Memory IPC: Copy, Move, BorrowRead, BorrowWrite, reply-bound revocation,
    cancellation, server-death recovery & Done \\
6 & CQ subsystem normalization: operation IDs, detached submission,
    CQ\_WAIT/CQ\_WAKE, per-shard CQ partitioning, backlog batching, 32-byte
    completion records & Done \\
7 & Sitas endpoint/CQ backend: endpoint readiness binding, unified shard
    wait, ShardExecutor (budgeted polling), ShardParker (spin free),
    per-shard CQ rings & Done \\
8 & Userspace UART driver: delegated MMIO + IRQ, EL0 MMIO writes,
    interrupt-driven deferred reads, crash recovery & Done \\
9 & Virtio-net driver: PCI discovery, MSI-X and SMMU delegation, feature
    negotiation, MAC/link read, virtqueue setup, and EL0 transmit submission.
    Smoltcp 0.13 adapter + TCP/IP service binary compile, and the frame
    demultiplexer routes IPv4/ARP to the tcpip service. &
    Two-guest TCG suite \\
10 & Lock-free device interrupt delivery (deferred wake), idempotent scheduler
    wakes, sustained-window load rebalancing & Done \\
11 & Userspace startup ABI: ELF \texttt{\_start}, \texttt{entry!} macro, typed
    \texttt{Context}, fixed-width versioned launch header, explicit startup
    input & Done \\
\hline
\end{longtable}

\section{Success criteria}

\begin{enumerate}
    \item Two EL0 domains communicate through a bounded endpoint --- \textbf{Verified.}
    \item Deferred async call while a shard continues other tasks --- \textbf{Verified.}
    \item Moved memory object --- sender locked out until return --- \textbf{Verified.}
    \item Lending enforced by mappings, including death cleanup --- \textbf{Verified.}
    \item Single \texttt{CQ\_WAIT} for endpoint + completions + device interrupts --- \textbf{Verified.}
    \item Terminal CQ results survive ring overflow (non-lossy backlog) --- \textbf{Verified.}
    \item Coalesced remote shard wakes (idempotent scheduler fix) --- \textbf{Verified.}
    \item Userspace UART driver with only delegated MMIO + IRQ --- \textbf{Verified.}
    \item Restart invalidates stale connections + device reset + op reconciliation --- \textbf{Verified.}
    \item Virtio-net data path with batched packet buffers --- \textbf{Runtime-validated on macOS QEMU TCG.}
    \item Fixed-width stable ABI representations (32-byte CQ entries) --- \textbf{Implemented; not formally audited.}
    \item Capability misuse and cancellation race adversarial tests --- \textbf{Partial; six IPC error-path tests.}
\end{enumerate}

\section{Verified scope}

\begin{itemize}
    \item \textbf{AArch64:} boots and the synchronous self-tests pass. The EL0
        service suite starts one shared node name service; Raft peers, echo,
        the service manager, NVMe/object store, UART, and their clients receive
        delegated connections to that registry. The deferred boot tests
        exercise registration, waitable lookup, restart, and live upgrade
        through \texttt{crt0/cmain}; the harness may print its synchronous
        success banner before those deferred tests finish. The sitas
        \texttt{basic\_kv} binary runs with 2 shards on per-shard CQ rings.

    \item \textbf{x86-64:} boots with \texttt{-cpu max}, core self-tests pass.
        Ring-3 userspace is implemented but not yet exercised.

    \item \textbf{Completion subsystem:} submit, complete, poll, wait, cancel,
        close, \texttt{submit\_worker} (exit-observer), CQ ring, CQ ring
        integrated with AS, non-lossy CQ overflow backlog, per-shard CQ
        partitioning, 32-byte entries.

    \item \textbf{Endpoint IPC:} bounded server queues, scalar messages,
        connection delegation with rights attenuation, memory-object attachments
        (move/lend/copy), deferred replies, reply-time connection delegation.

    \item \textbf{Device capabilities:} MmioRegion and Interrupt object types,
        user-accessible Device-nGnRnE page mapping, GICv3 SPI enable/route/mask,
        IRQ-to-CQ coalesced delivery (lock-free deferred wake).

    \item \textbf{Userspace drivers:} reference UART driver (PL011) with
        delegated MMIO + IRQ, interrupt-driven deferred reads, crash-test
        restart with device reset and operation reconciliation.

    \item \textbf{Syscall dispatch:} AArch64 SVC dispatch uses the shared typed
        \texttt{SyscallNumber} enum. Caller identity comes from trusted thread and
        trap context, never a userspace parameter. Explicit authenticated IPC
        receive operations add the sender's exact address-space generation,
        signed-artifact principal, and roles without changing the legacy
        nine-register receive ABI.

    \item \textbf{Authorization:} the local name service hosts a default-deny
        policy engine. Separate roles protect policy changes and service
        publication. Connection issuance is attenuated and generation-fenced;
        an administrator can read its bounded audit stream. Policy and audit
        persistence, distributed replication, and hard revocation remain
        future work.

    \item \textbf{Scheduler:} RoundRobin, quantum timer, block/yield/wake,
        idempotent wakes, deferred reaper, per-LP thread pinning.

    \item \textbf{Sitas integration:} sitas-core compiles for aarch64-none.
        sitas-charlotte links against the kernel's syscall ABI. Per-shard CQ
        rings mapped by the loader contract. ShardParker retires busy-waiting.

    \item \textbf{Protocol crates:} \texttt{charlotte-protocol-net},
        \texttt{charlotte-protocol-msg}.

    \item \textbf{Networking:} smoltcp 0.13 adapter and TCP/IP service runtime-validated over the frame demultiplexer; reliable-message fragmentation to 64 KiB.

    \item \textbf{Raft:} two-node election over local endpoint IPC, persistent
        single-voter restart, discovery-led join through joint consensus, and a
        two-guest replicated DNS catalog. Restart-safe admission has focused host
        tests and a persistent-storage boot-test manifest.

    \item \textbf{CI:} GitHub Actions builds both architectures with
        \texttt{-D warnings}, QEMU boot gate on AArch64.
\end{itemize}

\section{Known limitations}

\begin{itemize}
    \item \textbf{HVF does not honor hardware ASIDs.} Apple's
        Hypervisor.framework strips the ASID bits from \texttt{TTBR0\_EL1} on
        read and does not use them to isolate user translations, so the kernel's
        ASID-based TLB fast path is unsafe there. Under \texttt{hvf\_compat},
        \texttt{switch\_ctx} flushes the whole TLB on every context switch
        (\texttt{tlbi vmalle1is}) and the SVC handler attributes syscalls from a
        per-LP tracked logical ASID rather than a \texttt{ttbr0\_el1} readback.
        See \emph{AArch64 Port Status} for the full diagnosis. HVF also exposes
        no guest SMMU for the protected-DMA NVMe path. Its compatibility boot
        therefore runs 15 non-storage tests; the full 18-test suite and
        persistent-storage workflows require TCG.
    \item \textbf{x86-64 ring-3 userspace:} SYSCALL entry trampoline exists;
        the EL0 test and real-EL0 user thread are AArch64-only.
    \item \textbf{Virtio-net runtime:} initialization and transmit submission
        are runtime-validated on macOS under QEMU TCG. HVF remains unsuitable
        for direct EL0 device MMIO. The receive/transmit data path, frame
        demultiplexing, reliable messages, distributed name service, and the
        TCP/IP service are validated; TX reclamation and batching remain
        incomplete.
    \item \textbf{General process loader:} the ELF loader is fixed-address for
        reference service images, not a general loader with argv/env/auxv-style
        setup.
    \item \textbf{Resource accounting:} per-domain quotas are not yet
        implemented; capability-table capacity is fixed.
    \item \textbf{Cross-machine Raft boundary:} two TCG guests share a QEMU stream LAN;
        the reliable-message frame transport and remote Raft peer routing are
        connected and validated (\texttt{--dns-test}).
    \item \textbf{Observability HTTP:} the httpd keyhole aggregates the
        \texttt{observe} snapshot and per-service status into a JSON report
        served over the TCP/IP service (\texttt{--http-test}).
    \item \textbf{IOMMU/SMMU integration:} coherent Arm SMMUv3 discovery via
        ACPI IORT, per-requester DMA-domain capabilities, memory pinning,
        IOVA map/unmap, synchronous invalidation, fault reporting, teardown,
        and NVMe integration are implemented and QEMU-tested. Non-coherent
        systems, ATS/PRI, multiple SMMUs, and other IOMMU architectures remain.
    \item \textbf{Upgrade:} the EL0 service manager performs handoff and invokes
        the capability-checked replacement spawn path. It can
        fetch a replacement ELF from a well-known NVMe object, pass it as a
        memory capability, and fall back to the embedded recovery image.
        Structural ELF validation, immutable kernel snapshotting, and Ed25519
        CLS2 signature and logical-name checks are implemented. Rollback and key
        rotation policy, multiple state objects in the reference manager, and
        manager-owned teardown remain future work.
    \item \textbf{Resource bounds:} endpoint queues and rings are bounded, but
        the non-lossy CQ overflow backlog has no hard memory bound.
\end{itemize}

\endinput
